\textbf{Implantación y separación del recinto.}
La Figura~\ref{fig:plano-recinto-proyectado} presenta un módulo técnico nuevo de 10,0 por 7,0 m interiores y 3,0 m de altura libre. Estas dimensiones son condiciones proyectadas de obra; no describen el inmueble existente. El Anexo A4.4 desarrolla cargas estructurales, holguras, rutas y emplazamiento exterior \parencite[RT-06.01/03/06, pp.~14--15]{pucv-transversal}.

\begin{figure}[htbp]
\centering
\setstretch{1}
\begin{tikzpicture}[x=1.35cm,y=1.35cm,
  zona/.style={draw=SynNavyLight,fill=SynSurface,thick},
  texto/.style={font=\fontsize{9.5}{11.5}\selectfont,align=center},
  cota/.style={{Latex}-{Latex},thin,draw=SynNavyLight}]
  \draw[zona] (0,0) rectangle (10,7);
  \draw[zona] (0,0) rectangle (2.5,3.8);
  \draw[zona] (0,3.8) rectangle (2.5,7);
  \draw[zona] (7.5,0) rectangle (10,3.8);
  \draw[zona] (7.5,3.8) rectangle (10,7);
  \draw[zona] (2.5,1.5) rectangle (7.5,7);
  \draw[dashed,draw=SynNavyLight] (5,1.5)--(5,4.1);
  \draw[dashed,draw=SynNavyLight] (2.5,5)--(7.5,5);
  \node[texto,text width=29mm] at (1.25,5.4) {Baterías A\\2 cadenas\\40 bloques/cadena\\Aire separado de TI};
  \node[texto,text width=29mm] at (8.75,5.4) {Baterías B\\2 cadenas\\40 bloques/cadena\\Aire separado de TI};
  \node[texto,text width=28mm] at (1.25,1.9) {UPS A\\93T60\\Distribución A};
  \node[texto,text width=28mm] at (8.75,1.9) {UPS B\\93T60\\Distribución B};
  \node[texto,text width=28mm] at (3.75,2.9) {Cómputo\\24 U\\8 U ocupadas};
  \node[texto,text width=28mm] at (6.25,2.9) {Comunicaciones\\12 U\\10 U ocupadas};
  \draw[draw=SynNavyLight,fill=white,thick] (3.4,5.95) rectangle (4.6,6.65);
  \draw[draw=SynNavyLight,fill=white,thick] (5.4,5.95) rectangle (6.6,6.65);
  \node[texto] at (4,6.3) {Clima A};
  \node[texto] at (6,6.3) {Clima B};
  \draw[-{Latex},thick,draw=SynBlue] (4,5.95)--(4,4.2);
  \draw[-{Latex},thick,draw=SynBlue] (6,5.95)--(6,4.2);
  \draw[-{Latex},thick,dashed,draw=SynNavyLight] (3,4.2) |- (3.4,6.3);
  \draw[-{Latex},thick,dashed,draw=SynNavyLight] (7,4.2) |- (6.6,6.3);
  \node[texto,fill=SynSurface,text width=60mm] at (5,4.75) {Recinto común: 27,5 m$^2$\\TI y unidades interiores; sin muro entre zonas};
  \node[texto,fill=SynSurface,text width=55mm] at (5,3.9) {Suministro continuo / retorno discontinuo};
  \node[texto,text width=61mm] at (5,.75) {Circulación, acceso independiente\\frente útil $\geq1,20$ m};
  \draw[cota] (0,7.3)--(10,7.3);
  \node[texto,fill=white,inner sep=1pt] at (5,7.3) {10,0 m interiores};
  \draw[cota] (-.3,0)--(-.3,7);
  \node[texto,rotate=90,fill=white,inner sep=1pt] at (-.3,3.5) {7,0 m interiores};
  \node[texto,draw=SynNavyLight,fill=SynSurface,text width=65mm,anchor=north west,inner sep=2mm] at (0,-.45)
    {\textbf{Exterior A/B segregado}\\Cada grupo: zona $\geq6,0\times4,0$ m.\\Depósito propio: 1.500 L nominales.\\Volumen útil $\geq1.190$ L; contención neta $\geq1.650$ L.};
  \node[texto,draw=SynNavyLight,fill=SynSurface,text width=55mm,anchor=north east,inner sep=2mm] at (10,-.45)
    {\textbf{Trabajo y respaldo separados}\\Área de trabajo externa $\geq3,0\times2,0$ m.\\Stock local apagado en espacio $\geq2,0\times1,5$ m.\\Bancada y ensayos fuera de sala.};
\end{tikzpicture}
\caption{Plano de distribución proyectada del recinto técnico y zonas asociadas}
\label{fig:plano-recinto-proyectado}
\FuenteFigura{Elaboración de Synaptix: cotas y condiciones de obra nueva según \parencite[RT-06.03, p.~14]{pucv-transversal} y área de instalación según \parencite[secc.~2.2.6, p.~21]{p9-mext-manual-2024}. Las divisiones discontinuas delimitan usos dentro del recinto común; las flechas representan circulación de aire. El plano ejecutivo desarrolla holguras, seguridad de refrigerante y tratamiento de los recintos energéticos separados.}
\end{figure}


La superficie de 70 m$^2$ se distribuye en $2(2,5\times7)=35$ m$^2$ energéticos, $17,5+10=27,5$ m$^2$ comunes para TI/comunicaciones y climatización, y 7,5 m$^2$ de circulación. Generación y combustible permanecen fuera. Se mantienen dos gabinetes AR3104 de 24 U, uno exclusivo de servidores y otro de comunicaciones, con ocupaciones de proyecto de 8 y 10 U. Las rutas A/B, el acceso, la extracción de bancos y los respaldos tienen separación propia; la elevación y los cálculos de soporte se consultan en A4.4 y el inventario en T-11.

La implantación debe comprobar geometría, anclaje, evacuación, refrigerante, clasificación de gas y ubicación de los armarios de incendio efectivos. La superficie o el dibujo no acreditan cabida final ni conformidad instalada. El CLIENTE ejecuta las obras asignadas y Synaptix entrega y coordina los planos civil, estructural, unifilar, ambiental y de recepción, conservando la estrategia de continuidad de 4.2.1.
